\begin{tikzpicture}[font=\small,x=1cm,y=1cm]
\fill[gris](3.3,-.6)ellipse(.55 and 1);\draw[azul,thick](3.3,-.6)ellipse(.55 and 1);
\fill[gris](9,0)ellipse(2 and 2.2);\draw[azul,thick](9,0)ellipse(2 and 2.2);
\node[punto,label=above:{Foco}]at(0,0){};
\draw[rayo](0,0)--(2.86,0);\draw[azul!45,dashed](2.86,0)--(11.2,0);
\foreach \x/\n in {2.86/1,3.74/2,7/3,11/4}{\node[circle,fill=ocre,inner sep=2pt]at(\x,0){};\node[font=\footnotesize,above]at(\x,.1){$t_\n$};}
\draw[rayo](0,0)--(7.7316,1.701);
\node[circle,fill=turquesa,inner sep=2pt]at(7.7316,1.701){};
\node[nota,align=center]at(3.3,-1.7){Contorno anterior\\p. ej., extremidad};
\node[nota]at(9,-2.6){Segundo contorno cerrado};
\node[proceso,text width=13cm,anchor=north west]at(0,-3.2){\textbf{Sobre el rayo horizontal:} seleccionar $t_1=\min\{t_k>0\}$.\\La cara en $t_3$ puede mirar al foco y estar dentro del campo, pero queda oculta. Las caras de salida $t_2$ y $t_4$ tampoco reciben dosis de entrada en esta vía.};
\end{tikzpicture}
